\section{静态、仅保护与信息物理联合方法对比}\label{sec:rel-method-comparison}

\subsection{对比的公平性条件}
三种方法必须共享同一系统指纹、故障频率、切负荷定义、修复时间和年度小时数。只有保护与信息处理层逐级增加：
\begin{longtable}{@{}L{3.0cm}L{5.0cm}L{5.8cm}@{}}
\toprule
方法 & 事件逻辑 & 后果\\
\midrule
静态 FMEA & 故障发生即按未清除/全修复处理 & $S_u\tau_r$\\
仅保护 & 主/后备继电、断路器、配合与 FRT & 类条件三窗口\\
信息物理联合 & 仅保护事件树再条件于路径、QoS、共因、供电和电池 & 联合类条件三窗口\\
\bottomrule
\end{longtable}
若三行使用不同负荷倍率或不同故障率，差值不能解释为保护或信息系统贡献。

\subsection{静态 FMEA 基线}
对场景 $k$，静态基线把未清除切负荷保持到修复：
\begin{align}
 E_k^{static}&=S_{k,u}\tau_{k,r},\\
 \mathrm{EENS}^{static}&=\sum_k\lambda_kE_k^{static},\\
 \mathrm{LOLE}^{static}&=\sum_k\lambda_k\tau_{k,r}\mathbf1[S_{k,u}>\varepsilon],\\
 \mathrm{LOLF}^{static}&=\sum_k\lambda_k\mathbf1[S_{k,u}>\varepsilon].
\end{align}
它是保守基线，不执行保护逻辑，因此不应把其中的“未清除”解释为所有保护都拒动；它表示传统单持续时间 FMEA 的后果口径。

\subsection{仅保护模型}
仅保护模型删除全部信息元件、路径、环境和功能绑定，保留主后备轨迹及按需成功概率。其指标为
\begin{equation}
 \mathrm{EENS}^{P}=\sum_k\lambda_k\sum_{c_P}\pi(c_P\mid k)E(k,c_P).
\end{equation}
本地保护不依赖通信时，这一行反映真实保护快速清除的收益。若保护方案本来依赖通道，则为公平起见，应把通道作为保护链
固有成功概率或在联合行中绑定信息功能，并明确“仅保护”是理想通道反事实。

\subsection{信息物理联合模型}
联合行保留完整信息状态：
\begin{equation}
 \mathrm{EENS}^{ICP}=\sum_k\lambda_k\sum_{x^c,c_P}
 p(x^c\mid k)p(c_P\mid x^c,k)E(k,x^c,c_P).
\end{equation}
信息系统可以改变检测、主/后备跳闸通道、隔离和恢复，因而既改变保护类概率，也改变切换时长与残余切负荷。

\subsection{两个可解释差值}
平台报告
\begin{align}
 \Delta E_{cyber}&=\mathrm{EENS}^{ICP}-\mathrm{EENS}^{P},\\
 \Delta E_{protection}&=\mathrm{EENS}^{static}-\mathrm{EENS}^{P}.
\end{align}
$\Delta E_{cyber}$ 是相对理想信息条件的退化增量；$\Delta E_{protection}$ 是相对静态未清除基线的保护收益。二者不是
互斥 Shapley 分解，在非单调后果或误动存在时可以为负。

\subsection{单调性条件}
若对每个类满足
\begin{equation}
 E_{primary}\le E_{backup}\le E_{unclear},\qquad
 E_{auto}\le E_{manual}\le E_{failed},
\end{equation}
则降低信息可用率会把概率质量移向更差类，联合 EENS 不减；静态 FMEA 与信息可用率无关。该性质进入 E2E 回归：把信息
可用率从 $0.8$ 降至 $0.4$，静态值必须不变、联合值必须上升。若后果不满足单调性，测试应改为概率守恒和逐类解释，不能
强制错误的方向性。

\subsection{解析对照算例}
当前 HTTP 基准使用一项故障频率和三窗口后果，得到：
\begin{center}
\begin{tabular}{@{}lrrr@{}}
\toprule
方法 & EENS/(MWh/年) & LOLE/(h/年) & 解释\\
\midrule
静态 FMEA & 200.000000 & 20.000000 & 未清除后果保持到修复\\
仅保护 & 0.4007778 & 0.2000778 & 主/后备快速清除\\
信息物理联合 & 21.5217333 & 4.1600622 & 信息失效使部分事件进入人工/拒清类\\
\bottomrule
\end{tabular}
\end{center}
三行使用同一故障频率与切负荷数据。仅保护相对静态显著下降证明保护逻辑真正改变后果；联合相对仅保护上升证明信息路径
状态实际进入事件树，而非只显示一个可用率字段。

\subsection{为何不能直接与普通 FMEA 数值混比}
通用 \srcpath{run\_distribution\_fmea} 会从完整网络求解阶段切负荷，而三级对照入口消费用户提供的三窗口切负荷。
只有当三窗口后果由同一网络、同一故障和同一负荷倍率自动生成时，二者绝对值才可互比。当前 GUI 对照的目标是比较三种事件
逻辑，不宣称复现通用 FMEA 的每一网络后果。

\subsection{保护成功率敏感性}
主清除概率 $p_p=p_Rp_B$。在其余条件固定且后果单调时，主继电成功率导数为
\begin{equation}
 \frac{\partial\mathbb E[E]}{\partial p_R}
 =p_B(E_p-E_{fallback,B})+(1-p_B)(E_{fallback,R}-E_{fallback,B}),
\end{equation}
符号由后果差决定。提高成功率只有在主清除后果更优时才有正价值；这说明概率参数和区变异后果必须一起校准。

\subsection{配合裕度敏感性}
配合裕度不直接改变事件概率，而改变 \fld{coordination\_satisfied} 诊断。若应用策略规定裕度不足时闭锁自动重合或扩大隔离，
调用方应把诊断映射到不同后果类；当前对照保留实际清除结果，不把单纯裕度告警自动升级为停电。

\subsection{通信可用率敏感性}
单路径串联系统中，功能可用率近似为元件成功概率乘积；改善最薄弱串联元件的边际收益最大。冗余路径中，新增一条完全独立
路径的收益为现有全部路径失败概率乘新增路径成功率；若共享电源，收益显著降低。平台精确元件位枚举能够反映该差异。

\subsection{QoS 敏感性}
提高固有可用率不会挽救超过时延上限的路径。路径先经过确定性 QoS 门限，只有合格路径再参与概率枚举。因此优化通信网络时
应分别考察“路径是否满足实时契约”和“合格路径的可用率”，不能用单一九个九指标掩盖保护时延超限。

\subsection{共因敏感性}
两条名义冗余路径若共享共因组，其联合失效概率至少包含共因环境概率。即使每个独立元件可用率趋近一，功能不可用率也不会
低于该共因概率。联合对照中降低独立失效率但保持共因不变会出现收益饱和，这是合理现象而非枚举误差。

\subsection{供电和电池敏感性}
备用能量对持续时间呈阈值效应：$E_b/P_b\ge T_{outage}$ 时静态环境中节点保持可用，略低于阈值时整段视为失效；序贯模型
则在耗尽时刻发生状态切换。阈值附近应使用序贯事件队列或对停电持续时间积分，避免静态二值近似造成突跳。

\subsection{检测混淆敏感性}
提高总检出率但把概率从正确类别移到错误类别，可能扩大隔离区并增加 EENS。评估应同时报告 $P_D$、正确分类率和每类隔离
后果。只优化召回率会忽略错误分类的系统代价。

\subsection{恢复风险约束敏感性}
收紧最大不安全概率会删除高恢复量动作，短期 EENS 可能上升，但避免不安全带电和二次事故。可靠性和安全性不是同一目标；
平台把风险约束作为硬门槛，不能用较低 EENS 证明越过安全门槛的策略可接受。

\subsection{在线 DAE 与给定轨迹对比}
给定轨迹模式便于复现实验和批量统计，但轨迹不会因整定、换流器控制或支路跳闸反馈自动改变。在线 DAE 模式从系统状态生成
轨迹，并把清除动作反馈到第二次动态运行，能揭示 DER 限流导致的保护延时或 FRT 跳闸。两者的结果差异应解释为轨迹生成
保真度差异，而非随机误差。

\subsection{年度序贯与解析事件树对比}
解析事件树直接对互斥类求期望，无抽样误差，适合独立代表事件。序贯队列能处理事件重叠、电池跨事件状态和共因持续过程，
但有抽样误差。低事件率、每次事件前状态重置时，两者均值应一致；有记忆和重叠时，序贯结果是更合适的口径。

\subsection{验收矩阵}
\begingroup\footnotesize
\begin{longtable}{@{}L{4.1cm}L{5.1cm}L{5.1cm}@{}}
\toprule
验证对象 & 必须成立的关系 & 失败含义\\
\midrule
\endfirsthead
\toprule
验证对象 & 必须成立的关系 & 失败含义\\
\midrule
\endhead
\bottomrule
\endfoot
事件树概率 & 正概率类和为一 & 分支遗漏或重复\\
三级静态基线 & 不随信息可用率变化 & 比较基准被污染\\
联合后果单调算例 & 信息退化时 EENS 不减 & 功能绑定或后果选择错误\\
在线清除时间 & DAE 与事件树差不超过时间步 & 轨迹转换不一致\\
DER 轨迹身份 & 主/后备/未清除稳定 ID 一致 & 错位消费设备轨迹\\
重要抽样恒等 & $\kappa=1,w=1,N_{eff}=N$ & 似然比错误\\
精确枚举 & 状态概率和为一 & 退化元件或位映射错误\\
Birnbaum & 与有限差分一致 & 边际状态配对错误\\
保护误动分解 & EENS 分解残差为零 & 增量未进入总量或重复计数\\
区间并集 & LOLE 不超过各事件时长和 & 重叠停电重复计算\\
\end{longtable}
\endgroup

\subsection{结论使用规则}
若三级结果同时给出，应优先报告三行绝对值、两个差值和全部有效性标志。不能只展示“联合方法更差”而不说明这是对理想信息
的增量，也不能只展示“保护收益”而隐藏保护误动。所有百分比变化都应注明分母；静态基线接近零时不报告无意义的巨大百分比。
